\documentclass[11pt]{article}
\usepackage[margin=1.1in]{geometry}
\usepackage[cmex10]{amsmath}
\usepackage{amssymb,amsthm}
\usepackage[T1]{fontenc}
\usepackage[hidelinks,hypertexnames=false]{hyperref}
\allowdisplaybreaks[2]
\title{Balanced arrays at fixed row cap and arbitrary dimension}
\author{ProgramComputer}
\date{}
\newtheorem{theorem}{Theorem}

\DeclareFontShape{OMX}{cmex}{m}{n}{<-> cmex10}{}
\begin{document}
\maketitle

\begin{abstract}
For a fixed positive integer $r$, we study the polynomial whose
coefficient of $q^k$ counts labelled binary arrays with $n$ rows,
$d$ columns of sum $k$, and each row sum at most $r$.
The theorem gives a threshold in $n$ depending only on $r$
above which all its zeros are simple and strictly negative for every
$d\ge\max(2,r)$. The proof combines control of coefficient quotients
at the two ends with an integral estimate through the intervening
range of quotas.
\end{abstract}

\section{Statement and context}
The \href{https://www.symmetricfunctions.com/realRootedOpenProblems.htm\#conjBalancedArrays}{balanced-array conjecture of Per Alexandersson}
asks for real-rootedness throughout the parameter range.
Write $a_k$ for the number of the arrays in the abstract, and set
\[
  P_n^{(r,d)}(q)=\sum_{k=0}^{\lfloor nr/d\rfloor}a_kq^k.
\]
The following partial theorem gives a threshold uniform in dimension.
\begin{theorem}
For every integer $r\ge1$ there exists $N_r$ such that, for all integers
$n\ge N_r$ and $d\ge\max(2,r)$, the polynomial $P_n^{(r,d)}$ has exactly
$\lfloor nr/d\rfloor$ simple, strictly negative zeros.
\end{theorem}

The assertion is vacuous when $d>nr$. Thus each fixed cap has only
finitely many pairs $(n,d)$ left outside this assertion. The proof
gives an existential threshold and concerns real-rootedness and simplicity.
At the excluded one-column case $r=d=1$, the polynomial is $(1+q)^n$.
In particular the caps of any infinite sequence of distinct
counterexamples to real-rootedness must tend to infinity.

At every integer quota $0\le k\le D$, choose any row-size profile $(s_1,\dots,s_n)$ with $0\le s_j\le r$ and total $dk$, which is possible since $dk\le nr$. Assign successive blocks of lengths $s_j$ in the cyclic list of the $d$ column labels to the corresponding rows. No block repeats a label since $s_j\le d$, and each label occurs exactly $k$ times in the whole list of $dk$ positions. This realizes every such profile and gives $a_k>0$ throughout the degree range.

\section{The dimension range that remains}

A failure of the theorem would give a sequence with $n\to\infty$ at one fixed $r$. A bounded-dimensional subsequence is excluded by the fixed-pair eventual theorem when $r<d$ and by the unrestricted binomial multiplier theorem when $r=d\ge2$. We can therefore take $d\to\infty$. The rank-dependent dimension bound excludes $d\ge C_r\sqrt{n/\log n}$. On the remaining sequence $b=d^2/n\to0$. Write $T=nr/d$, $D=\lfloor T\rfloor$, so that $D/d\to\infty$. We prove the required signs for every sequence $d\to\infty$, $b\to0$ at fixed $r$. \href{https://github.com/ProgramComputer/open-math-research/blob/ba7ddfd7efb37cac33769f4fdfed749e1b5bf58a/topics/balanced-arrays/fixed-rank-eventual.pdf}{Fixed-pair theorem} \enspace \href{https://github.com/ProgramComputer/open-math-research/blob/ba7ddfd7efb37cac33769f4fdfed749e1b5bf58a/topics/balanced-arrays/finite-linear-reduction.md}{Unrestricted case} \enspace \href{https://github.com/ProgramComputer/open-math-research/tree/ba7ddfd7efb37cac33769f4fdfed749e1b5bf58a/topics/balanced-arrays/dense-regime/rank-dependent}{Rank-dependent theorem}

\section{Where the phases are attained}

For $w>0$ write $H(w)=\sum_{s=0}^r\binom ds w^s$ and let $S$ have these size weights, normalized by $H(w)$. Put $\mu=\mathbb E S$, $V=\mathrm{Var}(S)$, $y=n\mu/d$, $\phi=\pi/d$ and $\nu=(d-1)/2$. The tilt with mean quota $x$ is written $w_x$, so $y(w_x)=x$. Write $H_\phi=H(we^{i\phi})$. If $G(S)=S(d-S)/(d(d-1))$, put

\begin{equation*}
\lambda_\phi=\frac{\mathbb E[G(S)e^{i\phi S}]}{\mathbb E e^{i\phi S}},
\qquad \Psi(w)=n\arg H_\phi-\nu\arg\lambda_\phi.
\end{equation*}

All arguments use the branches obtained continuously from angle zero. For $d$ large, $r\phi<\pi/4$, and both positive-weight characteristic sums are in that angular sector. As $w$ increases from zero to infinity, $\Psi$ is continuous with limits $-\nu\phi$ and $nr\phi=T\pi$. The lower limit follows from the dominance of $S=1$ in the numerator defining $\lambda_\phi$, and the upper limit from $G(r)>0$ for $d>r$. Thus each internal level $k\pi$ has a first positive hitting point $v_k$, and $v_1<\cdots<v_{D-1}$ by continuity.

For the centered characteristic function, the imaginary part differs from zero by $O_r(V\phi^3)$ because the first moment vanishes and $\mathbb E|S-\mu|^3\le r V$. Its real part is $1+O(V\phi^2)$. This gives $\arg H_\phi=\mu\phi+O_r(V\phi^3)$. Moreover $|\arg\lambda_\phi|\le r\phi$. Bounded support gives $V\le C_r\min(\mu,r-\mu)$. Consequently, throughout the positive radius,

\begin{equation}
|\Psi(w)/\pi-y|\le C_r\big(1+\min(y,T-y)/d^2\big).
\tag{M1}
\end{equation}

At a crossing, (M1) implies $y=k+O(1)$ if $k=O(d)$ or $T-k=O(d)$: in the lower case bound the minimum by $y$, solve for $y$ and substitute, and at the upper end apply the same argument to $T-y$. If both $k/d$ and $(D-k)/d$ tend to infinity, so do $y/d$ and $(T-y)/d$, since a bounded subsequence on either side would contradict (M1).

At $k/d\to c\in(0,\infty)$, we have $\mu=O(b)$ at $v_k$. The law of $S$ concentrates at zero, and the law with weights multiplied by $G(S)$ concentrates at one. Indeed their respective exceptional probabilities are $O(\mu)$ from the expansion in $dw=O(\mu)$ and $G(s)/G(1)\le r$. For a bounded integer-valued law with probability $1-O(\mu)$ at $s_0$, its characteristic argument at $\phi$ is $s_0\phi+O_r(\mu\phi)$ by centering at $s_0$. Hence $\arg\lambda_\phi=\phi+O_r(\mu\phi)$. Also $nV\phi^3=O(k/d^2)=o(1)$ at the crossing. Therefore $y(v_k)=k+1/2+o(1)$. At $(D-k)/d\to c\in(0,\infty)$ we have $r-\mu=O(b)$, and both laws concentrate at $r$ with exceptional probability $O_r(r-\mu)$, since $G(s)/G(r)$ is bounded. This gives $\arg\lambda_\phi=O_r((r-\mu)\phi)$ and $nV\phi^3=O((T-k)/d^2)=o(1)$. Thus $y(v_k)=k+o(1)$ there. The tilt satisfies

\begin{equation*}
\frac{\mathrm d}{\mathrm dx}\big(d\log w_x\big)=\frac{b}{V_x},
\end{equation*}

and the endpoint expansions $V_x\sim dx/n$ at the lower end and $V_x\sim d(T-x)/n$ at the upper end give

\begin{equation}
d\log v_k=d\log w_k+\frac{d}{2k}+o(1)\quad(k/d\to c),
\qquad d\log v_k=d\log w_k+o(1)\quad((D-k)/d\to c).
\tag{M2}
\end{equation}

These expansions of $V_x$ hold on bounded translates of $k$ by expanding the finite size weights in $dw_x$ or its reciprocal. In particular $b/V_x\to1/c$ on either window.

\section{Alternation near each end}

Write $t_k=\sqrt{a_{k-1}/a_{k+1}}$ and $\Gamma_k=a_k^2/(a_{k-1}a_{k+1})$. The \href{https://github.com/ProgramComputer/open-math-research/blob/ba7ddfd7efb37cac33769f4fdfed749e1b5bf58a/topics/balanced-arrays/orientation-normalization.md}{normalized orientation inequality} implies strict log-concavity at every internal quota and $\Gamma_k\ge(1+1/k)^d$. On $k/d\to0$ this tends to infinity, including bounded $k$. On $(D-k)/d\to0$ the scalar deficit curvature and the uniform column-balance local limit prove $\Gamma_k\to\infty$. In either event $(-1)^kP(-t_k)>0$ eventually. To see this, divide by $a_kt_k^k$ and multiply by $(-1)^k$. The two immediate neighbors have magnitude $1/\sqrt{\Gamma_k}$, and each of the two remaining alternating tails begins with a positive term and has decreasing magnitudes by log-concavity. The resulting lower bound is $1-2/\sqrt{\Gamma_k}$. \hyperref[sec:upper]{Upper endpoint} \enspace \hyperref[sec:profiles]{Fixed-profile estimate}

At $k/d\to c\in(0,\infty)$ or $(D-k)/d\to c\in(0,\infty)$, the endpoint coefficient estimates give $\log\Gamma_{k+j}\to1/c$ for every fixed integer $j$. They also give $\log t_k=d\log v_k+o(1)$ by (M2). At any radius $x$ with $\log(x/t_k)=o(1)$, the ratios $a_{k+j}x^j/a_k$ tend to $\exp(-j^2/(2c))$. This follows by multiplying consecutive term ratios starting from the two neighbors at $t_k$. For any fixed large $L$, the outward consecutive ratios beyond distance $L$ are bounded by a constant $<1$ uniformly along the sequence, by log-concavity and the limit at $L$. The two sums beyond $L$ are therefore bounded by a constant times the boundary terms, which converge to $e^{-L^2/(2c)}$. Extending missing terms by zero and sending $L$ to infinity yields

\begin{equation*}
\frac{(-1)^kP(-x)}{a_kx^k}
\longrightarrow \sum_{j\in\mathbb Z}(-1)^j e^{-j^2/(2c)}
=\sqrt{2\pi c}\sum_{m\in\mathbb Z}e^{-2\pi^2c(m+1/2)^2}>0 .
\end{equation*}

The equality is Gaussian Poisson summation. This proves the sign at both $x=t_k$ and $x=v_k^d$. \hyperref[sec:lower]{Lower-end coefficients} \enspace \hyperref[sec:upper]{Upper-end coefficients}

\section{Alternation between the two ends}

Give the torus $|z_j|=1$, $\prod_j z_j=-1$ its normalized Haar measure. Write
\[R_w(z)=\sum_{|A|\le r} w^{|A|}\prod_{j\in A} z_j.\]
Then

\begin{equation*}
P(-w^d)=\int R_w(z)^n\,d\mathrm{Haar}(z).
\end{equation*}

Indeed all exponents must agree to survive averaging the first $d-1$ independent circle variables. The product constraint then gives $(-1)^k$ at quota $k$. At $w=v_k$ with $k/d,(D-k)/d\to\infty$, (M1) implies $n\mu/d^2,n(r-\mu)/d^2\to\infty$. Global localization gives an absolute error $o(|H_\phi|^nI_0)$ away from the two constant-phase neighborhoods. There $I_0=(2\pi)^{-(d-1)}d^{-1/2}(2\pi/(n\lambda_0))^{\nu}$ with $\lambda_0=\mathbb E G(S)$. For the common-arc tolerance $\eta$ in the localization estimate, choose the fixed centered hyperplane cube with coordinate deviations at most $\eta/2$, reducing its radius further if needed for the local expansion. Every pair of coordinates then differs by at most $\eta$, placing the entire cube in the corresponding principal chart. Its complement inside the short-arc cluster has negligible mass by the coordinate tail bound. The local transverse Gaussian evaluation on that cube gives the two conjugate terms, whose sum is

\begin{equation*}
2 |H_\phi|^n I_0 (\lambda_0/|\lambda_\phi|)^\nu
\big(\cos\Psi(w)+o(1)\big).
\end{equation*}

Here $|\lambda_\phi|/\lambda_0=1+O_r(\phi^2)$ since $\lambda_\phi/\lambda_0$ is the ratio of characteristic functions of two bounded-support laws at $\phi$. The determinant multiplier tends to one. The local term uses Euclidean hyperplane measure, whose Haar density is $(2\pi)^{-(d-1)}/\sqrt d$. Since $\Psi(v_k)=k\pi$, this yields $(-1)^kP(-v_k^d)>0$ eventually. \hyperref[sec:localization]{Global localization} \enspace \hyperref[sec:transverse]{Local Gaussian proof}

\section{One ordered set of test points}

Use the radius $t_k$ when $k\le d$ or $D-k\le d$, and $v_k^d$ otherwise. The two edge ranges are disjoint for large $d$ since $D/d\to\infty$. The geometric radii increase strictly by log-concavity, and the phase radii increase by their first-hitting definition. At either interface the neighboring geometric radii satisfy $\log(t_{k+1}/t_k)=(\log\Gamma_k+\log\Gamma_{k+1})/2\to1$ because the relevant distance divided by $d$ tends to one. Replacing one by its phase radius changes the logarithm by $o(1)$ by (M2) and the coefficient estimates. Thus all test radii are strictly ordered. Put zero at the left and take the right endpoint to positive infinity, with sign $(-1)^D$ there because the leading coefficient is positive.

If an interior sign fails along a subsequence, pass to a further subsequence on which $k/d$ and $(D-k)/d$ have limits in $[0,\infty]$. Their sum tends to infinity. A zero limit is in the geometric range, a positive finite limit is covered for either radius by the endpoint argument, and two infinite limits are in the phase range covered by localization. This excludes every failure. Alternating signs therefore give $D$ distinct negative real roots by the intermediate value theorem, exhausting the degree and proving simplicity.


\appendix
\clearpage
\section{Uniform local limit for balance at fixed row sizes}\label{sec:profiles}

Fix $r\ge1$. Let $d\to\infty$, take any finite list $s_1,\ldots,s_n\in\{0,\ldots,r\}$, and let $X_i$ be independent indicators of uniform $s_i$-subsets of $[d]$. Write $M=\sum_i s_i=dk$ for an integer $k$, and assume $M/d^2\to\infty$. Then, uniformly over those lists,

\begin{equation*}
\mathbb P\!\left(\sum_i X_i=k\mathbf 1\right)
=\sqrt d\,(2\pi\Lambda)^{-(d-1)/2}(1+o(1)),
\qquad
\Lambda=\sum_i\frac{s_i(d-s_i)}{d(d-1)}.
\end{equation*}

In particular $\Lambda\asymp M/d$. The local relative error is $O(d/\sqrt M+d^2/M)$ and the global relative tail is at most $d^2 C^d\exp(-c M/d)$, where constants depend only on the fixed cap.

\subsection{Quotient and normalization}

Let $f_s(\theta)=\mathbb E\exp(i\sum_{j\in S}\theta_j)$ for a uniform $s$-subset $S$. Pinned-coordinate Fourier inversion in angles $\theta_1,\ldots,\theta_{d-1}$ with $\theta_d=0$ gives the probability as the normalized integral of the function
\[\prod_i f_{s_i}(\theta)\exp(-ik\sum_j\theta_j).\]
The integrand is invariant under common phase shifts because $M=dk$. It is also periodic in each coordinate because $k$ is an integer.

Project orthogonally onto $V=\mathbf 1^\perp$ by $P=I-\mathbf 1\mathbf 1^\top/d$. The period lattice in $V$ has basis $2\pi Pe_1,\ldots,2\pi Pe_{d-1}$, as $Pe_d=-\sum_{j<d}Pe_j$. Their Gram determinant after removal of the $2\pi$ factors is $\det(I_{d-1}-\mathbf 1\mathbf 1^\top/d)=1/d$. The covolume of the projected pinned-coordinate cube is therefore $(2\pi)^{d-1}/\sqrt d$.

For any sufficiently small fixed $\eta>0$, the image of the cube $C_\eta=\{y\in V:\max_j|y_j|\le\eta\}$ is injective into the quotient. Indeed any period difference has pairwise coordinate differences that are integer multiples of $2\pi$, impossible for a nonzero period vector in $C_\eta-C_\eta$ if $4\eta<2\pi$. On this cube the centered integrand is $\prod_i f_{s_i}(y)$.

\subsection{Uniform local expansion}

For $y\in V$, put $L=\sum_{j\in S}y_j$. The inclusion probabilities are $h_j=(s)_j/(d)_j$. Its first moment vanishes, and direct expansion gives

\begin{equation*}
\mathbb E L^2=(h_1-h_2)\sum_j y_j^2,\qquad
\mathbb E L^3=(h_1-3h_2+2h_3)\sum_j y_j^3.
\end{equation*}

Here $\sum_{a\ne b}y_a^2y_b=-\sum_a y_a^3$ and the ordered distinct triple sum is $2\sum_a y_a^3$. The absolute fourth moment is at most $r^3(s/d)\sum_j y_j^4$ by the convexity inequality for a sum of $s$ terms. Choosing a fixed small cube makes $|f_s(y)-1|\le C(s/d)\sum_j y_j^2$ small, and squaring this bound is $O((s/d)\sum_j y_j^4)$ by Cauchy--Schwarz and fixed $r$. Hence

\begin{equation*}
\sum_i\log f_{s_i}(y)=-\frac{\Lambda}2\sum_j y_j^2
-\frac{i}{6}\gamma\sum_j y_j^3+E(y),\qquad
|\gamma|\le CM/d,\quad |E(y)|\le C(M/d)\sum_j y_j^4.
\end{equation*}

The constant $\gamma$ is real, and $h_1-h_2=s(d-s)/[d(d-1)]$. Zero sizes contribute nothing.

Let $Y_j=Z_j-\overline Z$ with independent $Z_j\sim N(0,\Lambda^{-1})$. This represents the Gaussian on $V$ with identity covariance on $V$ times $\Lambda^{-1}$. Put $q_{ab}=\Lambda^{-1}(\delta_{ab}-1/d)$. Wick gives $\mathbb E Y_a^3Y_b^3=9q_{aa}q_{bb}q_{ab}+6q_{ab}^3$. The first term cancels on summation since the diagonal entries agree and all row sums of $q$ vanish. Therefore $\|\sum_jY_j^3\|_2=O(\sqrt d\,\Lambda^{-3/2})$, and $\|\sum_jY_j^4\|_2=O(d\Lambda^{-2})$ by the eighth moment and triangle inequality. Writing $B=-i\gamma\sum_jY_j^3/6+E(Y)$ gives

\begin{equation*}
\|B\mathbf 1_{C_\eta}\|_2=O(d/\sqrt M+d^2/M).
\end{equation*}

For exponential integrability, bound $2(\mathrm{Re}B)_+$ by $C\Lambda\sum_jY_j^4$ on $C_\eta$. The event $|\overline Z|\le\eta$ has probability tending to one and is independent of $Y$. With that event and $Y\in C_\eta$, one has $|Z_j|\le2\eta$ and $\sum_jY_j^4\le16\sum_j Z_j^4$. Thus the required truncated exponential expectation is bounded by a constant times
$\bigl(\mathbb E[\mathbf 1_{|Z_1|\le2\eta}\exp(C'\Lambda Z_1^4)]\bigr)^d$.
Choose $\eta$ small enough that $C'\Lambda z^4\le\Lambda z^2/4$ on the interval. Using $e^t-1\le te^t$ and the fourth Gaussian moment with doubled variance bounds each factor by $1+O(1/\Lambda)$. Their product is bounded since $d/\Lambda=O(d^2/M)=o(1)$. Cauchy--Schwarz applied to $|e^B-1|\le|B|e^{(\mathrm{Re}B)_+}$ gives the local relative error. The missing Gaussian mass outside $C_\eta$ is at most $2d\exp(-c\eta^2\Lambda)$. The full Gaussian integral is $(2\pi/\Lambda)^{(d-1)/2}$. Multiplying by the normalized quotient density yields the stated constant.

\subsection{Global localization by one-element exchange}

Write $z_j=e^{i\theta_j}$, $\bar z=d^{-1}\sum_jz_j$, and $V_z=\sum_j|z_j-\bar z|^2$. For $s\ge1$, choose a uniform $(s-1)$-subset $A$ and then append $j$ uniformly from its complement of size $q=d-s+1$. This gives a uniform $s$-subset. By conditional Jensen, $1-|f_s|^2$ is at least the expected conditional variance of $z_j$. The variance on a fixed complement is the sum of ordered squared pair differences divided by $2q^2$. Averaging over the complement retains each pair with probability $q(q-1)/(d(d-1))$. Thus

\begin{equation*}
1-|f_s|^2\ge\frac{d(d-s)}{(d-1)(d-s+1)}\frac{V_z}{d}.
\end{equation*}

For sufficiently large $d$ and fixed $r$, this is at least $c_r(s/d)V_z$. The zero-size case is trivial. Consequently $|\prod_i f_{s_i}|\le e^{-c(M/d)V_z}$.

Choose an anchor $a$ minimizing $|z_a-\bar z|$. Then $\sum_j|z_j-z_a|^2=V_z+d|z_a-\bar z|^2\le2V_z$. On the quotient choose principal anchor angles $t_j=\arg(z_j/z_a)\in[-\pi,\pi]$ for $j\ne a$. These coordinates have the normalized Haar product measure, as can also be seen by redoing the pinned-coordinate parametrization with $a$ fixed. Chord length satisfies $|z_j-z_a|\ge 2|t_j|/\pi$. If every $|t_j|\le\eta/2$, the projection of the anchor angle vector with $t_a=0$ belongs to $C_\eta$. Therefore outside the local image at least one anchor coordinate exceeds $\eta/2$. Partition by the choice of closest anchor, then bound the integral by a sum of $d$ unrestricted pinned-coordinate integrals with that tail event. A union bound over the $d-1$ coordinates and one-dimensional Gaussian bounds give

\begin{equation*}
\left|\int_{\mathrm{outside}}\prod_i f_{s_i}(\theta)e^{-ik\sum_j\theta_j}\,d\mathrm{Haar}\right|
\le d^2 \left(\frac{C}{\sqrt{\Lambda}}\right)^{d-1}e^{-c\eta^2\Lambda}.
\end{equation*}

Its ratio to the main Gaussian contribution is at most $d^2 C_1^d e^{-c_1 M/d}=o(1)$ because $M/d^2\to\infty$. This proof uses no distribution or tail hypothesis for the row sizes beyond the fixed cap.

\section{Lower-edge coefficient comparison at fixed row cap}\label{sec:lower}

Fix an integer $r\ge1$. Let $n,d\to\infty$ with $b=d^2/n\to0$
and take integer quotas $k$ with $k/d\to c\in(0,\infty)$. Let $a_i$
count labeled binary $n\times d$ arrays with common column sum $i$
and row cap $r$. Define
\begin{equation*}
\Gamma_i=\frac{a_i^2}{a_{i-1}a_{i+1}},\qquad
t_i=\sqrt{\frac{a_{i-1}}{a_{i+1}}},\qquad
H(w)=\sum_{s=0}^r\binom ds w^s .
\end{equation*}
For real $u$ in a fixed bounded window about $k$, let $w_u>0$ be
the scalar capped-binomial tilt whose mean is $du/n$. Then for
every fixed integer $j$,
\begin{equation}
\log\Gamma_{k+j}\longrightarrow\frac1c,\qquad
\log t_k=d\log w_k+\frac{d}{2k}+o(1).
\tag{L1}
\end{equation}
All integer quotas involved are interior for large $n$, since
$\lfloor nr/d\rfloor/d\to\infty$.

\subsection{The row law at vanishing mean}

Write $\delta=d/n$, $\mu_u=\delta u$, and let $S$ have mass
$\binom ds w_u^s/H(w_u)$ for $0\le s\le r$. The tilt is unique
since the mean increases from zero to $r$ and its derivative
in $\log w$ is the positive variance. In the stated window,
$\mu_u\asymp b\to0$, whereas $n\mu_u\asymp d^2\to\infty$.
Writing $\lambda=dw_u$, the ratio of the masses at 1 and 0
is $\lambda$. Markov's inequality gives
$\mathbb P(S=0)\ge1-\mu_u$ and $\mathbb P(S=1)\le\mu_u$,
so $\lambda=O(\mu_u)$. The bound $\binom ds/d^s\le1/s!$ gives
\begin{equation*}
H(w_u)=1+\lambda+O(\lambda^2),\qquad
\mu_u=\lambda+O(\lambda^2),\qquad
\mathbb E[S(S-1)]=O(\mu_u^2).
\end{equation*}
In particular $\mathbb P(S=0)=1+O(\mu_u)$ and
$\mathbb P(S=1)\sim\mu_u$. With $V_u=\mathrm{Var}(S)$ we have
\begin{equation}
V_u=\mu_u+O(\mu_u^2),\qquad
\mathbb E|S-\mu_u|^3=O(\mu_u),\qquad
\left|\frac{dV_u}{d\mu_u}\right|=O(1).
\tag{L2}
\end{equation}
The last bound follows from the exponential-family identity
$dV/d\mu=\mathbb E(S-\mu)^3/V$. All constants are uniform on
the window.

At an integer quota $i$, let $T_m$ be the sum of $m$ independent
copies of the law at $w_i$, and put $M=di=n\mu_i$. Write $L_i$ for
$\mathbb P(T_n=M)$. Then
\begin{equation}
L_i=(2\pi nV_i)^{-1/2}(1+o(1)),\qquad
\sup_{v\in\mathbb Z}\mathbb P(T_{n-1}=v)=O((n\mu_i)^{-1/2}).
\tag{L3}
\end{equation}
Indeed for $\phi(t)=\mathbb E e^{it(S-\mu_i)}$ the masses at 0
and 1 in the cosine expansion of $1-|\phi|^2$ give
$|\phi(t)|\le e^{-c_0\mu_i t^2}$ on $[-\pi,\pi]$. Inversion and
absolute values yield the second estimate. At the exact mean
inversion uses $\phi(t)^n$ without an additional phase. With
$t=v/\sqrt{nV_i}$, the third-moment bound in (L2) gives
$\phi(t)^n\to e^{-v^2/2}$ uniformly for $v$ in bounded sets,
and the Fourier envelope yields Gaussian domination after
rescaling. This proves (L3) uniformly for bounded shifts of $k$.

\subsection{Summing the fixed-margin correction}

For integer $i$ in the window put
\begin{equation*}
A_i=[w^{di}]H(w)^n,\qquad K_i=\frac{\binom ni^d}{\binom{nd}{di}},
\qquad p_i=i/n.
\end{equation*}
For a row profile $\mathbf s$ with sum $M=di$ and $s_j\le r$,
the maximum row and column degrees are at most $r$ and $i$.
The Greenhill--McKay--Wang hypothesis
$ri/M^{2/3}=O(d^{-1/3})$ holds, and $M\to\infty$.
The normalized row variance is
$R/[n\mu_i(1-p_i)]$ for $R=\sum_j(s_j-\mu_i)^2$.
It is uniformly bounded since $\sum_j s_j^2\le rM$ and $p_i\to0$.
The normalized column variance is zero. Thus Corollary 5.2,
under Theorem 1.3, gives uniformly for every feasible row
profile the count
\begin{equation}
K_i\prod_j\binom d{s_j}
\exp\left(-\frac12+\frac{R}{2n\mu_i(1-p_i)}
                       +O(d^{-1/3})\right).
\tag{L4}
\end{equation}
Every profile of total $di$ under the cap is feasible. Assign successive blocks of lengths $s_j$ in the cyclic list of column labels to the rows. Since $s_j\le d$, the labels within each row are distinct, and each label is used exactly $i$ times across all $di$ positions.

No near-regularity or aspect ratio assumption is needed. \href{https://doi.org/10.1016/j.jcta.2005.03.005}{Greenhill--McKay--Wang, Theorem 1.3 and Corollary 5.2} / \href{https://web.maths.unsw.edu.au/~csg/papers/irregular-final.pdf#page=33}{Authors' text, Corollary 5.2}

The weights $\prod_j\binom d{s_j}/A_i$ at total $M$ are the
law of the $n$ tilted sizes conditioned on $T_n=M$.
With $U=\sum_j S_j(S_j-1)$, the exact identity on this event is
\begin{equation*}
R=M(1-\mu_i)+U,\qquad 0\le U\le(r-1)M.
\end{equation*}
Here the seemingly constant exponent $-1/2$ in (L4) cancels
the linear contribution of the squares. To control $U$ after
conditioning, use (L3):
\begin{equation*}
\begin{aligned}
\frac{\mathbb E[U\mid T_n=M]}M
&=\frac nM\sum_{s=0}^r s(s-1)\mathbb P(S=s)
  \frac{\mathbb P(T_{n-1}=M-s)}{\mathbb P(T_n=M)}\\
&=O(\mu_i).
\end{aligned}
\end{equation*}
The ratio of probabilities is uniformly bounded, even if the
number of rows grows arbitrarily fast relative to $d$. Now the
correction exponent without the $O(d^{-1/3})$ part equals
\begin{equation*}
\frac{p_i-\mu_i}{2(1-p_i)}+\frac{U/M}{2(1-p_i)}.
\end{equation*}
The first term is $O(\mu_i+p_i)$ and the second is bounded.
Since $|e^x-1|\le C|x|$ on a fixed bounded interval, its
conditional exponential average is $1+O(\mu_i+p_i)$. Combining
with the uniform error in (L4) gives
\begin{equation}
a_i=K_i A_i\bigl(1+O(\mu_i+p_i+d^{-1/3})\bigr).
\tag{L5}
\end{equation}
For $r=1$ the term $U$ is identically zero, so the same
calculation applies without any comparison of second factorial moments.

\subsection{Exact differences of the binomial normalization}

Set $\Delta_i=\log K_i-\log K_{i-1}$. The adjacent binomial
ratio and the product formula for the ratio at a distance $d$
give
\begin{equation*}
\Delta_i=\sum_{h=1}^d
\left[
  \log\frac{n-i+1}{n-i+h/d}
 +\log\frac{i-1+h/d}{i}
\right].
\end{equation*}
Writing $s_h=1-h/d$, the first logarithm is
$-\log(1-s_h/(n-i+1))$ and the second is
$\log(1-s_h/i)$. Since $\sum_h s_h=(d-1)/2$, Taylor expansion
with a uniform quadratic error gives
\begin{equation*}
\Delta_i=-\frac{d-1}{2i}
+O\left(\frac{d}{n-i+1}+\frac{d}{i^2}\right).
\end{equation*}
Thus $(\log K_{k-1}-\log K_{k+1})/2=d/(2k)+o(1)$.
Also differences of the same logarithms with $i$ incremented
by one are $O((n-i)^{-2}+i^{-2})$ termwise. Consequently
$2\log K_i-\log K_{i-1}-\log K_{i+1}=O(d/i^2+d/(n-i)^2)=o(1)$.
These estimates use exact products, so no Stirling remainder
is needed.

\subsection{Entropy differences and the answer}

At real quota $u$ write $E(u)=n\log H(w_u)-du\log w_u$.
At integers the exact tilted coefficient is $A_i=e^{E(i)}L_i$.
Differentiation gives
\begin{equation*}
E'(u)=-d\log w_u,\qquad E''(u)=-b/V_u,\qquad
E'''(u)=\frac{b\delta}{V_u^2}\frac{dV_u}{d\mu_u}=O(d^{-1}).
\end{equation*}
The last estimate uses (L2), $\mu_u\asymp b$, and $\delta=b/d$.
In the whole fixed window $b/V_u\to1/c$. Therefore, for fixed $j$,
the centered second difference of $E$ at $k+j$ tends to $1/c$
by the triangular integral. The centered first difference
$(E(k-1)-E(k+1))/2$ is $d\log w_k+O(d^{-1})$ by the bound on
the third derivative. By (L2) the ratios of the neighboring
variances tend to one, so (L3) makes each of the required
differences of $\log L_i$ equal to $o(1)$. Combining these
facts with (L5) and the binomial normalization calculation
proves (L1).

\section{Coefficients at the upper end for a fixed row cap}\label{sec:upper}

Fix a positive integer $r$, let $d\to\infty$, and put $b=d^2/n\to0$, $T=nr/d$ and $D=\lfloor T\rfloor$. Write $a_i$ for the number of labeled binary arrays with common column sum $i$ and row sums at most $r$. The functions $H(w)=\sum_{s=0}^r\binom ds w^s$ and $A_i=[w^{di}]H(w)^n$ count rows and arrays with only the total prescribed. At real quota $0<x<T$ choose $w_x$ by $\mu_x=\mathbb E S=dx/n$ for the size law with weights $\binom ds w_x^s$. Its variance $V_x>0$ is the derivative of its mean with respect to $\log w_x$, so the tilt exists uniquely. Put $h_x=r-\mu_x$.

For internal integers $k$ set $\Gamma_k=a_k^2/(a_{k-1}a_{k+1})$ and $t_k=\sqrt{a_{k-1}/a_{k+1}}$. If $(D-k)/d\to c\in(0,\infty)$, then, for every fixed integer $j$,

\begin{equation*}
\log\Gamma_{k+j}\to1/c,\qquad \log t_k=d\log w_k+o(1).
\end{equation*}

If $(D-k)/d\to0$, then $\Gamma_k\to\infty$.

\subsection{The column-balance probability}

For $0\le D-i=O(d)$, set $M=di$, $e=nr-M=d(T-i)$, and $m=d-1$. Thus $e=O(d^2)$ and $M\sim nr$. Condition on any row-size profile $s_1,\ldots,s_n$ of that total, with $0\le s_j\le r$. The uniform fixed-profile local limit gives probability $\sqrt d\,(2\pi\Lambda)^{-m/2}(1+o(1))$ of column balance for independent uniform supports of these sizes, uniformly even at $e=0$. \hyperref[sec:profiles]{Fixed-profile estimate}

Here $\Lambda=\sum_j s_j(d-s_j)/(d(d-1))$ and

\begin{equation*}
0\le\Lambda-\Lambda_*(i)=
\frac{\sum_j s_j(r-s_j)}{d(d-1)}\le\frac{re}{d(d-1)},
\qquad \Lambda_*(i)=\frac{i(d-r)}{d-1}.
\end{equation*}

Since $d\log(\Lambda/\Lambda_*)=O(e/M)=o(1)$, averaging over profiles gives $a_i=A_i F_i(1+o(1))$ for $F_i=\sqrt d\,(2\pi\Lambda_*(i))^{-m/2}$. The adjacent differences of $\log F_i$ are $O(d/i)=O(b)=o(1)$, since $i\sim nr/d$. This proves that the required local logarithmic differences may be taken on $A_i$, including when one neighbor is the last coefficient.

\subsection{Positive distance proportional to dimension}

If $(D-k)/d\to c>0$ is finite, then $h_x\sim cb$ uniformly for $|x-k|\le L$ at any fixed $L$. For $u=dw_x$, the size weights are proportional to $c_s u^s$ with $c_s=\binom ds/d^s$ bounded above and bounded away from zero for $0\le s\le r$ and $d\ge2r$. Expanding the law of the deficit $J=r-S$ in $1/u$ gives $V_x=h_x+O(h_x^2)$, $\mathbb P(J=0)=1-O(h_x)$ and $\mathbb P(J=1)\sim h_x$. The third absolute centered moment is $O(h_x)$ by bounded support, and $|dV_x/d\mu_x|\le C_r$ by the centered third-moment identity.

Write $E(x)=n\log H(w_x)-dx\log w_x$. Then $E'(x)=-d\log w_x$, $E''(x)=-b/V_x\to-1/c$ uniformly on these windows and $|E'''(x)|=O(1/d)$. At each integer $i$ there, $A_i=e^{E(i)}(2\pi nV_i)^{-1/2}(1+o(1))$. Indeed the adjacent masses at $J=0,1$ give $|\mathbb E e^{itJ}|\le e^{-c_r h_i t^2}$ throughout $[-\pi,\pi]$. After rescaling $t$ by $1/\sqrt{nV_i}$, the third-moment Taylor bound and $nh_i\asymp d^2\to\infty$ give pointwise Gaussian convergence in Fourier inversion at the integer mean. The preceding Gaussian bound justifies integrating the limit. Ratios of neighboring variances tend to one. The central second difference of $E$ is the triangular average of $-E''$, and its signed central first difference $(E(k-1)-E(k+1))/2$ is $-E'(k)+O(1/d)$. This proves both limits.

\subsection{Distances sublinear in dimension}

Here $d\le e=d(T-k)=o(d^2)$ at an internal quota. The scalar deficit counts $U_e=[z^e](\sum_{j=0}^r\binom d{r-j}z^j)^n$ satisfy log-concavity relative to $\binom{n+e-1}e$. For completeness, denote the positive row weights by $B_j$ and the $(n-1)$-row totals by $C_m$. The $C_m$ have interval support and are log-concave by convolution of log-concave sequences. The law on $j$ proportional to $(j+1)B_j C_{e-j}$ shifts upwards in likelihood ratio order when $e$ increases, by log-concavity of $C$. This holds at the support boundaries by extending $C_m$ by zero. Under that law the expectation of $B_{j+1}/B_j$ (zero at $j=r$) is $(e+1)U_{e+1}/((n+e)U_e)$, by differentiating the row polynomial power for the numerator and by symmetry of the $n$ summands for the denominator. The ratio $B_{j+1}/B_j$ decreases, hence this normalized forward ratio decreases. \href{https://github.com/ProgramComputer/open-math-research/blob/8bf2220688ce5fb1e4ac745daa655bb3c3813fbf/topics/balanced-arrays/deficit-curvature/scalar.md}{Original scalar derivation}

Comparing the binomial normalizers at $e-d,e,e+d$, and integrating the one-step differences of the logarithms, gives

\begin{equation*}
\log\frac{A_k^2}{A_{k-1}A_{k+1}}
\ge\sum_{j=1}^d\int_0^d\frac{n-1}{(e-d+j+t)(n+e-d+j-1+t)}\,dt
\ge\frac{d^2(n-1)}{(e+d)(n+e+d)}.
\end{equation*}

Since $e=o(d^2)$ and $d^2=o(n)$, the right side tends to infinity. The column-balance comparison transfers this divergence to $\log\Gamma_k$.

\section{Local transverse Gaussian estimate for fixed row cap}\label{sec:transverse}

Fix $r\ge1$, let $d\to\infty$, take $u>0$, and put $v=ue^{i\tau}/d$. Use the row function
$R_\tau(y)=\sum_{|S|\le r}v^{|S|}\exp(i\sum_{i\in S}y_i)$ on the Euclidean hyperplane $V=\{y\in\mathbb R^d:\sum_i y_i=0\}$. Put $H=R_\tau(0)$ and $\lambda=\min(u,1)$. All constants below can depend on fixed $r$.

For some fixed $\eta,\tau_0>0$, uniformly on $\max_i|y_i|\le\eta$, $|\tau|\le\tau_0$ and all $u>0$, the logarithm through zero satisfies

\begin{equation*}
\log(R_\tau(y)/H)=-\tfrac12 a\sum_i y_i^2-\tfrac i6 c\sum_i y_i^3+E_4(y),\qquad
|E_4(y)|\le C\lambda d^{-1}\sum_i y_i^4.
\end{equation*}

Here $a=a_1-a_2$ and $c=a_1-3a_2+2a_3$, where $a_j$ is the complex normalized weight of a subset containing any fixed $j$-set. In particular $\mathrm{Re}(a)\asymp\lambda/d$, $|\mathrm{Im}(a)|\le C\lambda|\tau|/d$, $|c|\le C\lambda/d$, and $|\mathrm{Im}(c)|\le C\lambda|\tau|/d$.

The local integral over the stated cube in $V$, with Euclidean measure, obeys

\begin{equation*}
\int_{V,\ \max_i|y_i|\le\eta} R_\tau(y)^n\,dy
=H^n\left(\frac{2\pi}{na}\right)^{(d-1)/2}
\left[1+O\left(\frac{d}{\sqrt{n\lambda}}+\frac{d^2}{n\lambda}\right)\right]
\end{equation*}

as long as $n\lambda/d^2\to\infty$ and $d\tau^2=O(1)$. Powers are continued from positive $a$. This includes the two principal phases $\tau=\pm\pi/d$ when $u$ stays in a positive compact and $n\ge c d^2\log n$. The assertion concerns this integral only. A global estimate and phase comparison are needed to identify the dominant contributions to the polynomial.

\subsection{Expansion}

Give each subset $S$ the normalized weight $v^{|S|}/H$ and write $L_S=\sum_{i\in S}y_i$. The positivity of the real radial coefficients and the bounded support of the sizes give $|H|\ge\tfrac12 H|_{\tau=0}$ for small enough $\tau_0$. Absolute complex weights are at most twice their $\tau=0$ counterparts. The real size distribution has $\mathbb P(|S|\ge1)\asymp\lambda$, uniformly in $u>0$, because the nonempty contribution is a polynomial in $u$ of degree $r$ with coefficients bounded above and below for large $d$. Its moments of positive order up to any fixed order are of size $\asymp\lambda$. This also controls derivatives of inclusion weights in $\tau$: $|a_j|+|\partial_\tau a_j|\le C\lambda/d^j$ for $1\le j\le3$. At zero phase $a_1\asymp\lambda/d$, $a_2=O(\lambda/d^2)$, and $a,c$ are real. These facts imply the asserted bounds for small $\tau_0$.

Weighted moments of $L_S$ are computed by exchangeability. The mean vanishes on $V$. The second moment is $(a_1-a_2)\sum_i y_i^2$. In the cubic, $\sum_{i\ne j}y_i^2 y_j=-\sum_i y_i^3$ and $\sum_{i,j,k\ {\rm distinct}}y_i y_j y_k=2\sum_i y_i^3$, giving the cubic coefficient. Weighted absolute fourth moments are at most $C\lambda d^{-1}\sum_i y_i^4$ by $|L_S|^4\le r^3\sum_{i\in S}|y_i|^4$. Taylor's bound for $e^{iL_S}$ gives the expansion through degree three. The logarithm exists in the stated cube because $|R_\tau/H-1|\le C\lambda d^{-1}\sum_i y_i^2\le C\eta^2$ using the vanishing first moment. The square of this bound is $O(\lambda d^{-1}\sum_i y_i^4)$ by Cauchy--Schwarz and $\lambda\le1$. The real part of the cubic term is at most $C\lambda|\tau|d^{-1}|\sum_i y_i^3|$.

\subsection{Integrated error}

Set $\alpha=\mathrm{Re}(a)$ and $\sigma^2=(n\alpha)^{-1}\asymp d/(n\lambda)$. A centered Gaussian $Y$ on $V$ of covariance $\sigma^2 I_V$ has coordinates $Y_i=Z_i-\overline Z$, with $Z_i$ independent $N(0,\sigma^2)$. Write $C_\eta$ for the cube and $B(y)=-inc\sum_i y_i^3/6+nE_4(y)$. The real Gaussian normalization divided by the absolute complex Gaussian normalization is $(|a|/\alpha)^{(d-1)/2}\le\exp(O(d\tau^2))$. The Gaussian mass outside the cube is at most $2d\exp(-c_0\eta^2 n\lambda/d)$, negligible compared with either displayed error scale.

For the unweighted $L^2$ norm of the cubic, the coordinate covariances are $q_{ij}=\sigma^2(\delta_{ij}-1/d)$. Wick's formula gives $\mathbb E Y_i^3Y_j^3=9q_{ii}q_{jj}q_{ij}+6q_{ij}^3$. Since all diagonal entries agree and $\sum_{ij}q_{ij}=0$, the linear terms cancel and $\|\sum_iY_i^3\|_2=O(\sqrt d\,\sigma^3)$. Also $\|\sum_iY_i^4\|_2\le C d\sigma^4$ by the eighth moment and the triangle inequality. Therefore

\begin{equation*}
\|B(Y)\mathbf 1_{C_\eta}\|_2
=O(d/\sqrt{n\lambda}+d^2/(n\lambda)).
\end{equation*}

To control the exponential in $|e^B-1|\le |B| e^{(\mathrm{Re} B)_+}$, bound $2(\mathrm{Re} B)_+$ by $C(n\lambda/d)(|\tau|\sum_i|y_i|^3+\sum_i y_i^4)$. The event $|\overline Z|\le\eta$ has probability tending to one independently of $Y$. On its intersection with $C_\eta$, one has $|Z_i|\le2\eta$ and $\sum_i|Y_i|^j\le 2^j\sum_i|Z_i|^j$ by the triangle inequality for powers and Jensen's inequality. Thus the expectation of the exponential over $C_\eta$ is at most a bounded factor times the product of $d$ identical one-dimensional truncated expectations. If $X\sim N(0,\sigma^2)$ and $L(x)=C'\sigma^{-2}(|\tau||x|^3+x^4)$, the single-factor bound is

\begin{equation*}
\mathbb E[\mathbf 1_{|X|\le2\eta}e^{L(X)}]
\le 1+C''(|\tau|\sigma+\sigma^2).
\end{equation*}

Indeed, after choosing small enough fixed $\eta$, $L(x)\le x^2/(4\sigma^2)$ on the interval. Use $e^{L(x)}-1\le L(x)e^{L(x)}$ there and integrate against the Gaussian, with variance at most doubled. The product is $\exp(O(|\tau|d^{3/2}/\sqrt{n\lambda}+d^2/(n\lambda)))$, bounded under the hypotheses. Cauchy--Schwarz now gives the claimed error after integrating the pointwise exponential difference and comparing normalizations.

\subsection{Moving radius}

The proof holds uniformly for $u\to0$ provided $nu/d^2\to\infty$, and for arbitrary $u\to\infty$ under $n/d^2\to\infty$. In the latter case the transverse variance stays of order $1/d$, even though the size distribution concentrates at $r$ and its radial variance can vanish. Consequently the local dimension estimates persist while a separate global argument in the radial/diagonal phase can lose localization. The proof gives no uniform vanishing-relative-error conclusion in the small-$u$ range $nu=O(d^2)$.

\section{Torus localization for fixed cap and a moving row tilt}\label{sec:localization}

2026-11-26 / Private research note

Fix $r\ge1$ and let $d\to\infty$, with $t=u/d>0$ and $u$ allowed to vary. Put

\begin{equation*}
H(w)=\sum_{s=0}^r\binom ds w^s,\qquad
R(\mathbf z)=\sum_{|A|\le r}t^{|A|}\prod_{i\in A}z_i,\qquad
H_\tau=H(te^{i\tau}),\qquad \phi=\pi/d .
\end{equation*}

Use normalized Haar measure on the torus $|z_i|=1$, $\prod_i z_i=-1$. Under the real row tilt $\mathbb P(A)=t^{|A|}/H(t)$, let $S=|A|$, $\mu=\mathbb E S$ and $V=\mathrm{Var}(S)$. Suppose that

\begin{equation*}
n\mu/d^2\longrightarrow\infty,\qquad
n(r-\mu)/d^2\longrightarrow\infty.
\end{equation*}

For distinct $i,j$ put $\lambda_0=\mathbb P(i\in A)-\mathbb P(i,j\in A)$. Then $\lambda_0\asymp_r\mu/d$. Use the positive reference saddle scale $|H_\phi|^n I_0$, where

\begin{equation*}
I_0=\frac{1}{(2\pi)^{d-1}\sqrt d}
\left(\frac{2\pi}{n\lambda_0}\right)^{(d-1)/2}.
\end{equation*}

The absolute integral off the two clusters of common phases $\pm\phi$ is $o(|H_\phi|^n I_0)$. The same applies within those clusters off the cutoff $\max_i|p_i|\le M\sqrt{d\log d/(n\mu)}$ for sufficiently large fixed $M$, where $p_i$ is the phase deviation from the common phase. All bounds below have constants depending only on $r$ and on fixed choices of small arcs. The reference scale is for a single saddle before adding the conjugate contribution.

\subsection{Row moments and diagonal gaps}

The size $s$ has weight $\binom ds (u/d)^s=a_su^s$, with $a_s$ bounded above and below by positive constants depending only on $r$ for $0\le s\le r$ and $d\ge2r$. For $0<u\le1$, the masses at zero and one are respectively $\asymp_r1$ and $\asymp_r u$. Here $\mu\asymp_r V\asymp_r u$ and $r-\mu$ is bounded below. For $u\ge1$, the masses at $r$ and $r-1$ are respectively $\asymp_r1$ and $\asymp_r u^{-1}$. Here $r-\mu\asymp_r V\asymp_r u^{-1}$ and $\mu$ is bounded below. The variance lower bounds follow by retaining the two indicated atoms in the independent-copy variance identity. The upper bounds follow by taking the squared distance from the dominant endpoint. Thus $V\asymp_r\min(\mu,r-\mu)$, and in particular $nV/d^2\to\infty$.

For any fixed $\eta_0>0$, the two respective adjacent masses give $|H_\tau/H(t)|\le e^{-cV}$ on $\eta_0\le|\tau|\le\pi$, by the squared-characteristic-function identity. Near zero, let $\chi(\tau)=\mathbb E e^{i(S-\mu)\tau}$. The centered moments of order $m\ge2$ in absolute value are bounded by $r^{m-2}V$. Hence $\chi=1+O(V\tau^2)$, $\chi'=O(V|\tau|)$, and $\chi''=-V+O(V|\tau|)$. On a sufficiently small fixed interval this gives $(\log|\chi|)''=-V+O(V|\tau|)\le -V/2$. By evenness,

\begin{equation*}
\log|H_\tau/H_\phi|\le -cV(\tau^2-\phi^2)
\quad\text{for }\phi\le|\tau|\le\eta_0,
\end{equation*}

when $\eta_0$ is chosen small enough. The same formulas also give $|H_\phi|/H(t)\ge \exp(-CV/d^2)$.

\subsection{Transverse bound and common-arc localization}

Let $\overline z=d^{-1}\sum_i z_i$. For $W=\prod_{i\in A}z_i$, the complex variance $\mathrm{Var}(W)=\mathbb E|W-\mathbb E W|^2$ is $1-|R/H(t)|^2$. Condition on $S=s\ge1$ and on the first $s-1$ indices of a uniform ordered rendering of $A$. The remaining index is uniform among $m=d-s+1$ available positions, and the product from the conditioned entries has modulus one. The resulting conditional variance is $1-|\overline z_{\mathrm{avail}}|^2$. Averaging over the conditioned indices gives

\begin{equation*}
\mathbb E(1-|\overline z_{\mathrm{avail}}|^2\mid S=s)
=\frac{d(m-1)}{m(d-1)}(1-|\overline z|^2).
\end{equation*}

Indeed, $1-|\overline z_{\mathrm{avail}}|^2=m^{-2}\sum_{i<j\text{ available}}|z_i-z_j|^2$ and each pair is available with probability $m(m-1)/(d(d-1))$. The prefactor is bounded below for $d\ge2r+1$, and $\mathbb P(S\ge1)\ge\mu/r$. The law of total variance therefore proves

\begin{equation}
|R/H(t)|\le \exp[-c\mu(1-|\overline z|^2)].
\tag{G1}
\end{equation}

Choose a representative coordinate $j$ minimizing $\sum_i|z_i-z_j|^2$. This minimum is at most the average $2d(1-|\overline z|^2)$. Fix a small angle $\eta>0$. Outside the common-arc region, for every representative $j$ some other coordinate $i$ satisfies $|\arg(z_i/z_j)|>\eta$. For the minimizing $j$, retain half of that chord penalty in (G1), giving a factor $e^{-c n\mu/d}$ after taking the $n$th power. Eliminate a coordinate different from $i,j$ in the torus parametrization. Conditional on the freely integrated pin $z_j$, the other $d-2$ coordinates in this chart have independent Haar measure and each contributes at most $C\sqrt{d/(n\mu)}$ from its chord factor. Summing over the at most $d(d-1)$ pairs, the absolute integral outside the common-arc region is bounded by

\begin{equation*}
H(t)^n d^2 e^{-c n\mu/d}(C d/(n\mu))^{(d-2)/2}.
\end{equation*}

Relative to $|H_\phi|^n I_0$, its logarithm is bounded above by $-c n\mu/d+O(d+\log d+\log(n\mu/d)+nV/d^2)$. This tends to minus infinity, since $n\mu/d\gg d$ and $V=O_r(\mu)$.

Every common-arc point can be lifted to angles differing by at most $2\eta$ using a representative. Their average $\tau$ is an odd multiple of $\pi/d$ by the product constraint. On reducing $\tau$ modulo $2\pi$ to $(-\pi,\pi]$, write $z_i=e^{i(\tau+p_i)}$ with $\sum_i p_i=0$ and $|p_i|\le2\eta$. This representation is unique within these bounds: if two representations coincide, differences of their componentwise integer angle shifts are zero since the pairwise differences of the deviations have range smaller than $2\pi$. The shifts thus coincide at every coordinate, and averaging shows that the canonical phases and deviations coincide. Each cluster uses measure $(2\pi)^{-(d-1)}dp_1\cdots dp_{d-1}$ with $p_d=-\sum_{i<d}p_i$. This follows by using the first $d-1$ circle variables, each with a unique lift in the short interval for fixed $\tau$. The uniqueness avoids a factor of $d$ in the later sum over small-angle clusters. On every such cluster, the pairwise chord bound and $\sum_i p_i=0$ give $1-|\overline z|^2\ge c d^{-1}\sum_i p_i^2$.

\subsection{Far common phases}

For $X_A=\sum_{i\in A}p_i$, conditional symmetry gives $\mathbb E(X_A\mid S)=0$ and

\begin{equation*}
|R/H(t)-H_\tau/H(t)|
=\left|\mathbb E\bigl[e^{i\tau S}(e^{iX_A}-1-iX_A)\bigr]\right|
\le C \frac{\mu}{d}\sum_i p_i^2.
\end{equation*}

Here $\mathbb E X_A^2\le r\mathbb E\sum_{i\in A}p_i^2=(r\mu/d)\sum_i p_i^2$. For $|\tau|\ge\eta_0$, write $B=(\mu/d)\sum_i p_i^2$. If $B\le c_1 V$ for sufficiently small $c_1>0$, the diagonal gap and this Taylor bound give $|R/H(t)|\le 1-c_2V\le e^{-c_3(V+B)}$. If $B>c_1V$, (G1) gives $|R/H(t)|\le e^{-cB}\le e^{-c_4(V+B)}$. Integrating the resulting transverse Gaussian over the zero-sum hyperplane and summing the at most $d$ clusters bounds their contribution relative to $|H_\phi|^n I_0$ by $\exp[-c nV+O(d+\log d+nV/d^2)]$. This is negligible.

\subsection{Small common phases and the sharp transverse volume}

Choose $\eta_0$ and then $\eta$ small enough in terms of $r$. For $|\tau|\le\eta_0$, the complex tilt with weights proportional to $e^{i\tau S}t^{|A|}$ is well-defined since $|H_\tau/H(t)|\ge\cos(r\eta_0)>0$. Denote expectations under it by $\mathbb E_\tau$. Direct differentiation of the weights shows that, for $m=2,3,4$, their $m$th moment of $X_A$ and its first two $\tau$ derivatives are at most $C(\mu/d)\sum_i|p_i|^m$ in modulus. In fact $\mathbb E|X_A|^m\le r^{m-1}(\mu/d)\sum_i|p_i|^m$, and the size factors from differentiation are uniformly bounded. Symmetry gives $\mathbb E_\tau X_A=0$ and $\mathbb E_\tau X_A^2=\lambda_\tau\sum_i p_i^2$, where $\lambda_\tau$ is the difference between the one- and two-coordinate inclusion weights. At zero,

\begin{equation*}
\frac{\mu}{d}\left(1-\frac{r-1}{d-1}\right)\le\lambda_0\le\frac{\mu}{d}.
\end{equation*}

The derivative bounds and conjugation give $\mathrm{Re}\lambda_\tau=\lambda_0+O(\mu\tau^2/d)$. The real part of the cubic term in $\mathbb E_\tau e^{iX_A}$ has size $O(|\tau|(\mu/d)\sum_i|p_i|^3)$ since it vanishes at zero. Taylor expansion through the cubic and the logarithm, using $|\tau||p_i|^3\le(\tau^2p_i^2+p_i^4)/2$, give

\begin{equation}
\log|R/H_\tau|
\le-\frac{\lambda_0}{2}(1-C\tau^2)\sum_i p_i^2
 +\frac{C\mu}{d}\sum_i p_i^4.
\tag{G2}
\end{equation}

The fourth-order Taylor remainder obeys the moment bound above. The logarithm error is bounded by $C(\mu d^{-1}\sum_i p_i^2)^2\le C(\mu/d)\sum_i p_i^4$, since $\mu\le r$. These Taylor expansions apply uniformly once the arcs are small enough.

Set $a=n\lambda_0(1-C\tau^2)/2\asymp n\mu/d$, $b=Cn\mu/d$ and $f(x)=e^{-ax^2+bx^4}\mathbf1_{\{|x|\le2\eta\}}$. The arc can be chosen so that the quartic term is at most half the quadratic on the support and $f$ decreases for positive $x$. The majorant in (G2) integrates in the $d-1$ free coordinates to at most $f^{*d}(0)/(2\pi)^{d-1}$. The one-dimensional mass is $Z=\sqrt{\pi/a}(1+O(1/a))$: after scaling $y=\sqrt a\,x$, the difference from the Gaussian has absolute integral at most $O(1/a)$ by the bound $(C/a)y^4e^{-y^2/2}$ plus a Gaussian cutoff tail.

The normalized convolution densities satisfy $\|(f/Z)^{*m}\|_\infty\le C\sqrt{a/m}$ for $m\ge2$. In the scaled variable the density tends in $L^1$ to a normalized Gaussian, with variance bounded below and fourth moment bounded above. The characteristic modulus is at most $e^{-c\xi^2}$ near zero by the independent-copy variance bound and fourth-order cosine Taylor bound. On every fixed interval away from zero it is bounded above by a constant smaller than one by the Gaussian limit. At larger frequencies it is at most $C/|\xi|$ by integration by parts, since the monotone halves give bounded total variation. Taking the split to large frequencies beyond $2C$, the integral of the $m$th power is $O(m^{-1/2})$ and Fourier inversion proves the assertion.

Consequently each small-angle cluster contributes at most

\begin{equation*}
C I_0 |H_\tau|^n \exp(Cd\tau^2+Cd^2/(n\mu)).
\end{equation*}

The first correction comes from replacing $\lambda_0$ by $\lambda_0(1-C\tau^2)$ and the second from $Z^d$. For a nonprincipal cluster $|\tau|\ge3\phi$, the first correction is absorbed in the diagonal gap because $nV\gg d$ and $\tau^2\le(9/8)(\tau^2-\phi^2)$. Summing over the odd phases gives, relative to $|H_\phi|^n I_0$, at most

\begin{equation*}
C e^{Cd^2/(n\mu)}
\sum_{j\ge1}\exp[-c(nV/d^2)((2j+1)^2-1)]
=O(e^{-c'nV/d^2})=o(1).
\end{equation*}

For either principal cluster, flag a coordinate with $|p_i|>v$. Integrate the normalized density of that coordinate over the tail, bounded by $C e^{-cav^2}$, and use the $(d-1)$-fold convolution sup bound for the remaining coordinates. Union over $d$ coordinates gives an extra factor $Cd e^{-cav^2}$. At $v=M\sqrt{d\log d/(n\mu)}$, this is $o(1)$ for large fixed $M$. This proves the localization and the cutoff.

\end{document}
